% !TeX root = ../main.tex
\section{Machine learning and decoding primer}
\label{sec:learning}

\subsection{Supervised learning in one page}

A \term{model} is a function with many adjustable numbers (\term{parameters} or
\term{weights}). \term{Training} adjusts them so that, on example inputs, the
model's outputs approach known correct answers (\term{labels}), by minimising a
\term{loss} that measures the disagreement. For transcription, the input is a
spectrogram and the labels say, for each frame and pitch, whether a note starts
or sounds there.

Data is split three ways. The \term{training set} is what the model learns
from. The \term{validation set} is used to make choices during development
(when to stop training, which settings to keep). The \term{test set} is used
only to report the final number; if any choice is made by looking at it, the
number becomes optimistic. A test set is \term{held out} when no part of it, nor
the same performers, appears in training. A result is \term{zero-shot} when the
model is tested on a dataset it never saw at all, which is the situation of a
user's song.

\begin{caution}[: overfitting and leakage]
A model that scores far better on its training players than on new ones has
\term{overfit}: it memorised instead of generalising. Our fingering tagger
scores 0.95 string accuracy on its training players and 0.57--0.69 on new
ones (Section~\ref{sec:fingering-app}). \term{Leakage} is subtler: if test
recordings share performers, songs or rendered tones with the training data,
the test score is inflated. The rule in this project is that GuitarSet players
03--05 are never used to fit anything, so they stay a clean test of fingering.
\end{caution}

\term{Fine-tuning} starts from a model already trained on one dataset and
continues training on another, so a small new dataset is enough. Our electric
guitar model is the GAPS classical-guitar model fine-tuned on electric guitar.
\term{Data augmentation} creates extra training examples by transforming the
ones we have; here, playing clean DI recordings through many amplifier
settings so the model sees the same notes in many tones.

\subsection{Network families mentioned in this document}

\begin{description}[style=nextline,leftmargin=1.2cm]
\item[CNN (convolutional neural network)] Slides small learned filters over the
  spectrogram image; good at local patterns such as a note's harmonic stack.
\item[RNN / GRU / LSTM (recurrent networks)] Read a sequence one step at a time
  and carry a memory; good at how notes evolve over frames.
\item[CRNN] A CNN followed by a recurrent layer. The GAPS guitar model (Kong et
  al.'s architecture) is a CRNN.
\item[Transformer] Relates every position of a sequence to every other one
  through \term{attention}~\cite{vaswani2017attention}. Used by BS-RoFormer,
  Beat This!, MT3-style transcribers and our fingering tagger.
\item[Encoder-decoder / sequence-to-sequence] An encoder reads the input; a
  decoder writes an output sequence one \term{token} at a time, each token
  conditioned on the ones before (\term{autoregressive} decoding). MT3,
  YourMT3+ and MuScriptor write notes as tokens such as ``time 1.23 s'',
  ``instrument: distorted guitar'', ``note on 52''. Decoding one token at a time
  is slow on a CPU.
\item[Diffusion model] Learns to turn random noise into a sample step by step;
  Noise2Fret uses one to generate tablature.
\end{description}

\subsection{From network output to notes}

Most note models output a \term{posteriorgram}: for each frame and each of the
88 piano pitches, a number between 0 and 1 that the model believes is the
probability of an onset (or of a sounding note) there. Turning it into a list of
notes requires \term{post-processing}: find peaks above a threshold in the
onset map, follow the note while the frame map stays above another threshold,
drop notes that are too short. These thresholds are not learned and change
results noticeably; in the app they live next to each model
(\file{transcribe/guitar_crnn.rs}, \file{transcribe/basic_pitch.rs}).

\begin{inapp}[: how models run]
All neural models used inside the app are exported to ONNX (a standard file
format for networks) and run in pure Rust with the \code{tract-onnx} crate
(version 0.23), on the CPU. Basic Pitch is compiled into the binary; the guitar
CRNN and Beat This! are optional downloads; the fingering tagger and the
electric CRNN are local files the user copies in. Source separators are the
exception: they run as external Python programs (MSST or Demucs) in their own
virtual environments, on the GPU when PyTorch finds one.
\end{inapp}

\subsection{Dynamic programming and the Viterbi algorithm}
\label{sec:viterbi}

Several stages must pick the best \emph{sequence} of choices: a string for each
note, a subdivision for each beat, a bar length for each bar. Choosing each
item on its own ignores how consecutive choices fit together (moving the hand
from fret 2 to fret 14 between two fast notes is bad even if each position is
fine alone). Trying every sequence is impossible: 20 notes with 3 positions each
give $3^{20} \approx 3.5$ billion sequences.

\begin{keyidea}[: Viterbi decoding]
Define a cost for each choice on its own (\term{local cost}) and a cost for each
pair of consecutive choices (\term{transition cost}). The best sequence ending in
choice $j$ at step $t$ is the best sequence ending in some choice $i$ at step
$t-1$, plus the transition $i\to j$, plus the local cost of $j$. So we only need
to keep, for each choice at each step, the cost of the best path that reaches it
and a pointer to where it came from. At the end we follow the pointers back.
The work is (steps)~$\times$~(choices)$^2$ instead of
(choices)$^{\text{steps}}$~\cite{viterbi1967error}.
\end{keyidea}

\begin{worked}[: two notes, three positions]
Note 1 can be played at fret 0, 5 or 10; note 2 at fret 2, 7 or 12. Local cost:
0.1 per fret (prefer low positions). Transition cost: 0.3 per fret of hand
movement. Best costs after note 1: $\{0, 0.5, 1.0\}$. For note 2 at fret 7:
from fret 0 costs $0 + 0.3\cdot 7 + 0.7 = 2.8$; from fret 5 costs $0.5 + 0.3\cdot
2 + 0.7 = 1.8$; from fret 10 costs $1.0 + 0.9 + 0.7 = 2.6$. Keep 1.8 with a
pointer to fret 5. Doing the same for frets 2 and 12 and taking the cheapest
end gives fret 0 then fret 2 (total $0 + 0.6 + 0.2 = 0.8$). If note 1 had been
forced high, the transition cost would have pulled note 2 high too: that is
the hand-position memory the method provides.
\end{worked}

A \term{hidden Markov model} (HMM) is the probabilistic version of the same
idea: costs are negative log-probabilities. A \term{dynamic Bayesian network}
(DBN) generalises it to several hidden variables (for beats: tempo and position
in the bar). Our fingering, our beat subdivision choice and our bar decoding are
all Viterbi searches with hand-set costs.
